\titre{}
\theme{numerique}
\auteur{Nathan Scheinmann}
\niveau{1M}
\source{musy}
\type{serie}
\piments{2}
\pts{}
\annee{2425}

\contenu{
\tcblower
Décrire les ensembles suivants par une condition d'appartenance (comme dans l'énoncé de l'activité 4).
\begin{center}
\begin{inlineumerate}
\item $\{\ldots\,;\,-3\,;\,-1\,;\,1\,;\,3\,;\,5\,;\,7\,;\,9\,;\,11\,;\,13\,;\,\ldots\}$
\item $\{0,;\,2\,;\,4\,;\,6\,;\,8\,;\,\ldots\}$
\item $\{1\,;\,4\,;\,9\,;\,25\,;\,\ldots\}$
%\item $\left(^*\right) \left\{0\,;\,\dfrac{1}{3}\,;\,\dfrac{1}{2}\,;\,\dfrac{3}{5}\,;\,\dfrac{2}{3}\,;\,\dfrac{5}{7}\,;\,\ldots\right\}$
\end{inlineumerate}
\end{center}
}
\correction{
\tcblower
\begin{tasks}(3)
		\task $\{2n+1 \mid n\in \Z\}$
		\task $\{2n \mid n\in \N\}$
		\task $\{n^2 \mid n\in \N^*\}$
	\end{tasks}
}

